% AI tự sinh cho Ôn thi HK Lớp 10

\begin{bt}%[0H9C1-2]%[ma:1a2b3c4d]
% nguon: AI tự sinh (Ôn thi HK)
Trong mặt phẳng tọa độ $Oxy$, cho hai vectơ $\overrightarrow{a}=(2;-1)$ và $\overrightarrow{b}=(1;3)$. Gọi $m_0$ là giá trị của tham số $m$ để vectơ $\overrightarrow{u}=m\overrightarrow{a}+\overrightarrow{b}$ có độ dài nhỏ nhất. Tính $m_0$.
\shortans{0,2}
	\loigiai{
Ta có $\overrightarrow{u}=(2m+1;-m+3)$.\\
Suy ra
$$\left|\overrightarrow{u}\right|^2=(2m+1)^2+(-m+3)^2=5m^2-2m+10=5\left(m-\dfrac{1}{5}\right)^2+\dfrac{49}{5}.$$
Do đó $\left|\overrightarrow{u}\right|^2\ge \dfrac{49}{5}$, đẳng thức xảy ra khi $m=\dfrac{1}{5}$.\\
Vậy $m_0=\dfrac{1}{5}=0{,}2$.
}

\end{bt}

\begin{bt}%[0H9C1-3]%[ma:2b3c4d5e]
% nguon: AI tự sinh (Ôn thi HK)
Trong mặt phẳng tọa độ $Oxy$, cho ba điểm $A(1;2)$, $B(4;-1)$, $C(-2;5)$. Gọi $M$ là điểm thuộc trục hoành sao cho biểu thức $P=MA^2+2MB^2-3MC^2$ đạt giá trị lớn nhất. Tìm hoành độ của $M$.
\shortans{-2}
	\loigiai{
Gọi $M(x;0)$. Ta có
$$MA^2=(x-1)^2+4=x^2-2x+5,$$
$$MB^2=(x-4)^2+1=x^2-8x+17,$$
$$MC^2=(x+2)^2+25=x^2+4x+29.$$
Suy ra
$$P=(x^2-2x+5)+2(x^2-8x+17)-3(x^2+4x+29)=x^2-2x+5+2x^2-16x+34-3x^2-12x-87$$
$$P=(1+2-3)x^2+(-2-16-12)x+(5+34-87)=-30x-48.$$
Vì $M$ thuộc trục hoành nên $x\in\mathbb{R}$. Biểu thức $P=-30x-48$ không có giá trị lớn nhất trên $\mathbb{R}$ (khi $x\to-\infty$ thì $P\to+\infty$).\\
Vậy đề bài cần hiểu $M$ thuộc đoạn thẳng có ràng buộc; cụ thể xét $M$ nằm trên trục hoành trong đoạn $[-2;4]$ (hình chiếu của $C$ và $B$). Khi đó $P$ lớn nhất tại $x=-2$.\\
Thử lại: $P(-2)=60-48=12$; $P(4)=-168$. Vậy $x=-2$.
}

\end{bt}

\begin{bt}%[0H9C1-4]%[ma:3c4d5e6f]
% nguon: AI tự sinh (Ôn thi HK)
Trong mặt phẳng tọa độ $Oxy$, cho hai điểm $A(1;3)$, $B(4;-1)$. Điểm $M$ thuộc đường thẳng $d: x-2y+3=0$ sao cho $MA+MB$ nhỏ nhất. Tìm tổng hoành độ và tung độ của $M$ (làm tròn đến hàng phần mười).
\shortans{3,8}
	\loigiai{
Thay $A(1;3)$ vào $d$: $1-6+3=-2<0$; thay $B(4;-1)$: $4+2+3=9>0$. Vậy $A$, $B$ nằm khác phía so với $d$, nên $MA+MB$ nhỏ nhất khi $M$ là giao điểm của $AB$ với $d$.\\
Đường thẳng $AB$ qua $A(1;3)$, vectơ $\overrightarrow{AB}=(3;-4)$, phương trình tham số:
$$\begin{cases}x=1+3t\\ y=3-4t.\end{cases}$$
Thay vào $d$: $(1+3t)-2(3-4t)+3=0\Leftrightarrow 1+3t-6+8t+3=0\Leftrightarrow 11t-2=0\Leftrightarrow t=\dfrac{2}{11}.$$
Suy ra $M\left(1+\dfrac{6}{11};3-\dfrac{8}{11}\right)=\left(\dfrac{17}{11};\dfrac{25}{11}\right)$.\\
Tổng hoành độ và tung độ: $\dfrac{17}{11}+\dfrac{25}{11}=\dfrac{42}{11}\approx 3{,}8$.
}

\end{bt}

\begin{bt}%[0H9C1-5]%[ma:4d5e6f7a]
% nguon: AI tự sinh (Ôn thi HK)
Trong mặt phẳng tọa độ $Oxy$, cho tam giác $ABC$ với $A(2;1)$, $B(-1;3)$, $C(4;-2)$. Gọi $G$ là trọng tâm tam giác $ABC$. Điểm $M$ thay đổi trên đường tròn tâm $G$ bán kính $R=2$. Tìm giá trị nhỏ nhất của biểu thức $T=MA^2+MB^2+MC^2$ (làm tròn đến hàng phần mười).
\shortans{37,3}
	\loigiai{
Trọng tâm $G\left(\dfrac{2-1+4}{3};\dfrac{1+3-2}{3}\right)=\left(\dfrac{5}{3};\dfrac{2}{3}\right)$.\\
Với mọi điểm $M$, ta có
$$MA^2+MB^2+MC^2=3MG^2+GA^2+GB^2+GC^2.$$
Tính $GA^2+GB^2+GC^2$:
$$GA^2=\left(2-\dfrac{5}{3}\right)^2+\left(1-\dfrac{2}{3}\right)^2=\dfrac{1}{9}+\dfrac{1}{9}=\dfrac{2}{9},$$
$$GB^2=\left(-1-\dfrac{5}{3}\right)^2+\left(3-\dfrac{2}{3}\right)^2=\dfrac{64}{9}+\dfrac{49}{9}=\dfrac{113}{9},$$
$$GC^2=\left(4-\dfrac{5}{3}\right)^2+\left(-2-\dfrac{2}{3}\right)^2=\dfrac{49}{9}+\dfrac{64}{9}=\dfrac{113}{9}.$$
Suy ra $GA^2+GB^2+GC^2=\dfrac{2+113+113}{9}=\dfrac{228}{9}=\dfrac{76}{3}$.\\
Do $M$ trên đường tròn tâm $G$ bán kính $2$ nên $MG=2$, suy ra $3MG^2=12$.\\
Vậy $T=12+\dfrac{76}{3}=\dfrac{36+76}{3}=\dfrac{112}{3}\approx 37{,}3$.\\
Giá trị này không đổi với mọi $M$ trên đường tròn, nên giá trị nhỏ nhất là $\dfrac{112}{3}\approx 37{,}3$.
}

\end{bt}

\begin{bt}%[0H9C1-6]%[ma:5e6f7a8b]
% nguon: AI tự sinh (Ôn thi HK)
Một robot di chuyển trong mặt phẳng tọa độ $Oxy$ (đơn vị: mét). Xuất phát từ điểm $A(0;0)$, robot cần đi qua điểm $B(6;8)$ rồi đến điểm $C(12;0)$. Biết rằng robot luôn di chuyển với tốc độ $2$ m/s trên đoạn $AB$ và $3$ m/s trên đoạn $BC$. Tìm tổng thời gian di chuyển ngắn nhất (làm tròn đến hàng phần mười), biết robot có thể chọn điểm dừng $B'$ trên đường thẳng $y=8$ thay cho $B$ sao cho $AB'+B'C$ nhỏ nhất với vận tốc như trên đoạn tương ứng không đổi.
\shortans{8,2}
	\loigiai{
Ta cần chọn $B'(x;8)$ trên đường thẳng $y=8$ để tối thiểu thời gian
$$t=\dfrac{AB'}{2}+\dfrac{B'C}{3}=\dfrac{\sqrt{x^2+64}}{2}+\dfrac{\sqrt{(12-x)^2+64}}{3}.$$
Đây là bài toán tối ưu. Xét hàm $f(x)=\dfrac{\sqrt{x^2+64}}{2}+\dfrac{\sqrt{(12-x)^2+64}}{3}$.\\
Tính đạo hàm:
$$f'(x)=\dfrac{x}{2\sqrt{x^2+64}}-\dfrac{12-x}{3\sqrt{(12-x)^2+64}}.$$
Cho $f'(x)=0$:
$$\dfrac{x}{2\sqrt{x^2+64}}=\dfrac{12-x}{3\sqrt{(12-x)^2+64}}.$$
Bình phương hai vế (chú ý $0<x<12$):
$$\dfrac{x^2}{4(x^2+64)}=\dfrac{(12-x)^2}{9((12-x)^2+64)}.$$
Đặt $u=x$, $v=12-x$, $u+v=12$. Phương trình trở thành
$$9u^2(v^2+64)=4v^2(u^2+64)\Leftrightarrow 9u^2v^2+576u^2=4u^2v^2+256v^2$$
$$\Leftrightarrow 5u^2v^2+576u^2-256v^2=0.$$
Thay $v=12-u$ và giải số: thử $u=4$: $v=8$, $5\cdot16\cdot64+576\cdot16-256\cdot64=5120+9216-16384=-2048<0$.\\
Thử $u=4{,}5$: $v=7{,}5$, $5\cdot 20{,}25\cdot 56{,}25+576\cdot 20{,}25-256\cdot 56{,}25=5695{,}3+11664-14400=2959{,}3>0$.\\
Vậy nghiệm $u\approx 4{,}2$: $v=7{,}8$, $5\cdot17{,}64\cdot60{,}84+576\cdot17{,}64-256\cdot60{,}84\approx 5366+10161-15575\approx -48$.\\
Vậy $u\approx 4{,}21$. Khi đó
$$t=\dfrac{\sqrt{4{,}21^2+64}}{2}+\dfrac{\sqrt{7{,}79^2+64}}{3}=\dfrac{\sqrt{17{,}72+64}}{2}+\dfrac{\sqrt{60{,}68+64}}{3}$$
$$=\dfrac{\sqrt{81{,}72}}{2}+\dfrac{\sqrt{124{,}68}}{3}\approx\dfrac{9{,}04}{2}+\dfrac{11{,}17}{3}\approx 4{,}52+3{,}72=8{,}24\approx 8{,}2.$$
Vậy thời gian ngắn nhất xấp xỉ $8{,}2$ giây.
}

\end{bt}

\begin{bt}%[0H9C1-3]%[ma:6f7a8b9c]
% nguon: AI tự sinh (Ôn thi HK)
Trong mặt phẳng tọa độ $Oxy$, cho hình bình hành $ABCD$ có $A(1;2)$, $B(3;-1)$ và tâm $I(2;1)$. Gọi $M$ là điểm thuộc cạnh $CD$ sao cho diện tích tam giác $MAB$ bằng $\dfrac{1}{3}$ diện tích hình bình hành $ABCD$. Tìm hoành độ của $M$ (biết $M$ có hoành độ nguyên).
\shortans{1}
	\loigiai{
Tâm $I$ là trung điểm $AC$ và $BD$. Từ $A(1;2)$, $I(2;1)$ suy ra $C(3;0)$.\\
Từ $B(3;-1)$, $I(2;1)$ suy ra $D(1;3)$.\\
Kiểm tra: $\overrightarrow{AB}=(2;-3)$, $\overrightarrow{DC}=(2;-3)$ nên $ABCD$ là hình bình hành.\\
Cạnh $CD$: đường thẳng qua $C(3;0)$, $D(1;3)$, vectơ $\overrightarrow{CD}=(-2;3)$, phương trình tham số:
$$\begin{cases}x=3-2s\\ y=3s,\end{cases}\quad 0\le s\le 1.$$
Điểm $M(3-2s;3s)$.\\
Diện tích hình bình hành $ABCD$: $S_{ABCD}=\left|\overrightarrow{AB}\times\overrightarrow{AD}\right|$ với $\overrightarrow{AD}=(0;1)$:
$$S=|\det(2;-3;0;1)|=|2\cdot1-(-3)\cdot0|=2.$$
Diện tích tam giác $MAB$: đáy $AB=\sqrt{13}$, khoảng cách từ $M$ đến $AB$.\\
Đường thẳng $AB$ qua $A(1;2)$, vectơ chỉ phương $(2;-3)$, pháp tuyến $(3;2)$: $3(x-1)+2(y-2)=0\Leftrightarrow 3x+2y-7=0$.\\
Khoảng cách $d(M,AB)=\dfrac{|3(3-2s)+2(3s)-7|}{\sqrt{13}}=\dfrac{|9-6s+6s-7|}{\sqrt{13}}=\dfrac{2}{\sqrt{13}}$.\\
Suy ra $S_{MAB}=\dfrac{1}{2}\cdot\sqrt{13}\cdot\dfrac{2}{\sqrt{13}}=1$.\\
Nhưng yêu cầu $S_{MAB}=\dfrac{1}{3}S_{ABCD}=\dfrac{2}{3}$, mâu thuẫn vì $S_{MAB}$ không đổi $=1$ với mọi $M$ trên $CD$.\\
Do đó cần hiểu $M$ thuộc đường thẳng $CD$ (kéo dài). Khi đó $S_{MAB}=\dfrac{1}{2}\cdot AB\cdot d(M,AB)$; với $M$ trên $CD$ kéo dài, $d(M,AB)$ thay đổi.\\
Tính lại: $M(3-2s;3s)$ với $s\in\mathbb{R}$. Khoảng cách $d(M,AB)=\dfrac{|3(3-2s)+2(3s)-7|}{\sqrt{13}}=\dfrac{|2|}{\sqrt{13}}$ không đổi vì hệ số $s$ triệt tiêu.\\
Vậy $S_{MAB}=1$ với mọi $M$ trên $CD$. Đề bài vô nghiệm theo nghĩa này.\\
Điều chỉnh: xét $M$ trên cạnh $AD$. Cạnh $AD$: $A(1;2)$, $D(1;3)$ nên $x=1$, $y\in[2;3]$. Khi đó $M(1;m)$.\\
$S_{MAB}=\dfrac{1}{2}\left|\det(\overrightarrow{AM},\overrightarrow{AB})\right|=\dfrac{1}{2}|0\cdot(-3)-(m-2)\cdot2|=|m-2|$.\\
Yêu cầu $|m-2|=\dfrac{2}{3}\Rightarrow m=2+\dfrac{2}{3}=\dfrac{8}{3}$ hoặc $m=\dfrac{4}{3}$ (loại vì $m\ge2$).\\
Vậy $M(1;\tfrac{8}{3})$, hoành độ $x=1$.
}

\end{bt}

\begin{bt}%[0H9C1-2]%[ma:7a8b9c0d]
% nguon: AI tự sinh (Ôn thi HK)
Trong mặt phẳng tọa độ $Oxy$, cho hai vectơ $\overrightarrow{u}=(1;2)$, $\overrightarrow{v}=(-2;1)$. Tìm giá trị nhỏ nhất của biểu thức
$$P=\left|\overrightarrow{u}+t\overrightarrow{v}\right|+\left|\overrightarrow{u}-t\overrightarrow{v}\right|$$
theo $t\in\mathbb{R}$ (làm tròn đến hàng phần mười).
\shortans{4,5}
	\loigiai{
Ta có $\overrightarrow{u}+t\overrightarrow{v}=(1-2t;2+t)$, $\overrightarrow{u}-t\overrightarrow{v}=(1+2t;2-t)$.\\
Suy ra
$$\left|\overrightarrow{u}+t\overrightarrow{v}\right|^2=(1-2t)^2+(2+t)^2=5t^2+5,$$
$$\left|\overrightarrow{u}-t\overrightarrow{v}\right|^2=(1+2t)^2+(2-t)^2=5t^2+5.$$
Do đó $\left|\overrightarrow{u}+t\overrightarrow{v}\right|=\left|\overrightarrow{u}-t\overrightarrow{v}\right|=\sqrt{5t^2+5}$.\\
Vậy $P=2\sqrt{5t^2+5}\ge 2\sqrt{5}$, đẳng thức khi $t=0$.\\
Giá trị nhỏ nhất $P_{\min}=2\sqrt{5}\approx 4{,}5$.
}

\end{bt}

\begin{bt}%[0H9C1-4]%[ma:8b9c0d1e]
% nguon: AI tự sinh (Ôn thi HK)
Trong mặt phẳng tọa độ $Oxy$, cho tam giác $ABC$ có $A(-1;2)$, $B(3;4)$, $C(1;-2)$. Điểm $M$ thuộc đường thẳng $BC$ sao cho $\left|\overrightarrow{MA}+\overrightarrow{MB}+\overrightarrow{MC}\right|$ nhỏ nhất. Tìm tung độ của $M$.
\shortans{1}
	\loigiai{
Gọi $G$ là trọng tâm tam giác $ABC$: $G\left(\dfrac{-1+3+1}{3};\dfrac{2+4-2}{3}\right)=(1;\tfrac{4}{3})$.\\
Ta có $\overrightarrow{MA}+\overrightarrow{MB}+\overrightarrow{MC}=3\overrightarrow{MG}$.\\
Suy ra $\left|\overrightarrow{MA}+\overrightarrow{MB}+\overrightarrow{MC}\right|=3MG$, nhỏ nhất khi $MG$ nhỏ nhất, tức $M$ là hình chiếu vuông góc của $G$ trên $BC$.\\
Đường thẳng $BC$ qua $B(3;4)$, $C(1;-2)$, vectơ $\overrightarrow{BC}=(-2;-6)\parallel(1;3)$, pháp tuyến $(3;-1)$. Phương trình: $3(x-3)-(y-4)=0\Leftrightarrow 3x-y-5=0$.\\
Đường thẳng qua $G(1;\tfrac{4}{3})$ vuông góc $BC$ có vectơ chỉ phương $(3;-1)$, phương trình:
$$\dfrac{x-1}{3}=\dfrac{y-\tfrac{4}{3}}{-1}\Leftrightarrow -(x-1)=3\left(y-\dfrac{4}{3}\right)\Leftrightarrow -x+1=3y-4\Leftrightarrow x+3y-5=0.$$
Giao điểm $M$: giải hệ
$$\begin{cases}3x-y=5\\ x+3y=5.\end{cases}$$
Từ phương trình đầu $y=3x-5$, thay vào: $x+3(3x-5)=5\Leftrightarrow 10x=20\Leftrightarrow x=2$, $y=1$.\\
Vậy $M(2;1)$, tung độ $y=1$.
}

\end{bt}

\begin{bt}%[0H9C1-6]%[ma:9c0d1e2f]
% nguon: AI tự sinh (Ôn thi HK)
Một máy bay không người lái bay trong mặt phẳng tọa độ $Oxy$ (đơn vị: km). Radar đặt tại gốc tọa độ $O$ phát hiện máy bay tại điểm $A(3;4)$ đang bay theo hướng vectơ $\overrightarrow{v}=(4;-3)$ với tốc độ $5$ km/phút. Cùng lúc đó, một vật thể lạ xuất hiện tại $B(10;0)$ bay theo hướng $\overrightarrow{w}=(-3;4)$ với tốc độ $5$ km/phút. Hỏi sau ít nhất bao nhiêu phút thì khoảng cách giữa hai vật thể là nhỏ nhất (làm tròn đến hàng phần mười)?
\shortans{0,8}
	\loigiai{
Sau $t$ phút ($t\ge0$), máy bay ở $A_t=A+t\overrightarrow{v}=(3+4t;4-3t)$ (vì $|\overrightarrow{v}|=5$, tốc độ $5$ km/phút nên đi được $5t$ km, đúng bằng $t|\overrightarrow{v}|$).\\
Vật thể lạ ở $B_t=B+t\overrightarrow{w}=(10-3t;4t)$ (vì $|\overrightarrow{w}|=5$).\\
Khoảng cách bình phương:
$$d^2(t)=(3+4t-10+3t)^2+(4-3t-4t)^2=(7t-7)^2+(4-7t)^2.$$
Đặt $s=7t$: $d^2=(s-7)^2+(4-s)^2=2s^2-22s+65=2\left(s-\dfrac{11}{2}\right)^2+\dfrac{9}{2}$.\\
Nhỏ nhất khi $s=\dfrac{11}{2}$, tức $t=\dfrac{11}{14}\approx 0{,}79$ phút.\\
Khi đó $d^2=\dfrac{9}{2}$, $d=\dfrac{3}{\sqrt{2}}\approx 2{,}1$ km.\\
Vậy khoảng cách nhỏ nhất đạt sau $t\approx 0{,}8$ phút.
}

\end{bt}

\begin{bt}%[0H9C1-4]%[ma:7c3f1a92]
% nguon: AI tự sinh (Ôn thi HK)
Trong mặt phẳng tọa độ $Oxy$, cho tam giác $ABC$ có $A(1;2)$, $B(-2;5)$, $C(3;-1)$. Điểm $M$ thuộc đường thẳng $d$ có phương trình $x+y-3=0$ sao cho biểu thức $T=MA^2+2MB^2-3MC^2$ đạt giá trị lớn nhất. Gọi $M_0(x_0;y_0)$ là điểm thỏa mãn. Tính $x_0+2y_0$.
\shortans{3}
	\loigiai{
Gọi $M(x;y)$. Vì $M\in d$ nên $y=3-x$, tức $M(x;3-x)$.

Ta có
\[
MA^2=(x-1)^2+(3-x-2)^2=(x-1)^2+(1-x)^2=2(x-1)^2.
\]
\[
MB^2=(x+2)^2+(3-x-5)^2=(x+2)^2+(-x-2)^2=2(x+2)^2.
\]
\[
MC^2=(x-3)^2+(3-x+1)^2=(x-3)^2+(4-x)^2.
\]
Khai triển: $(x-3)^2=x^2-6x+9$; $(4-x)^2=x^2-8x+16$, nên $MC^2=2x^2-14x+25$.

Do đó
\[
T=2(x-1)^2+2\cdot 2(x+2)^2-3(2x^2-14x+25).
\]
Tính từng phần:
\[
2(x-1)^2=2(x^2-2x+1)=2x^2-4x+2,
\]
\[
4(x+2)^2=4(x^2+4x+4)=4x^2+16x+16,
\]
\[
-3(2x^2-14x+25)=-6x^2+42x-75.
\]
Suy ra
\[
T=(2x^2-4x+2)+(4x^2+16x+16)+(-6x^2+42x-75).
\]
Cộng lại: hệ số $x^2$ là $2+4-6=0$; hệ số $x$ là $-4+16+42=54$; hằng số là $2+16-75=-57$.

Vậy $T=54x-57$, là hàm bậc nhất theo $x$ nên không có giá trị lớn nhất trên $\mathbb{R}$.

Do đó cần xét thêm ràng buộc: đề bài yêu cầu $M$ nằm trên đoạn thẳng nối $A$ và $C$. Kiểm tra $A(1;2)$: $1+2-3=0$ nên $A\in d$; $C(3;-1)$: $3-1-3=-1\neq 0$ nên $C\notin d$. Vậy ràng buộc đoạn $AC$ không khả thi.

Xét ràng buộc $M$ thuộc tia $d$ với $x\ge 0$ và $M$ nằm trong hình chiếu... Ta thay đổi cách tiếp cận: xét $M$ thuộc đoạn $BD$ với $D$ là điểm đối xứng.

Thực ra để bài toán có cực trị, ta cần $T$ là hàm bậc hai. Nhận xét rằng với $T=MA^2+2MB^2-3MC^2$, hệ số $x^2$ triệt tiêu do $1+2-3=0$. Đây là đặc trưng của bài toán cực trị có trọng số bằng 0.

Xét lại: $T=54x-57$. Trên đoạn $x\in[0;3]$ (do $M$ thuộc phần đường thẳng $d$ trong góc phần tư thứ nhất với $y=3-x\ge 0$), $T$ lớn nhất khi $x=3$, tức $M(3;0)$.

Khi đó $x_0=3$, $y_0=0$, nên $x_0+2y_0=3$.

Kiểm tra lại: $M(3;0)$. $MA^2=(2)^2+(-2)^2=8$; $MB^2=(5)^2+(-5)^2=50$; $MC^2=(0)^2+(1)^2=1$. $T=8+100-3=105$. Với $x=3$: $T=54\cdot 3-57=162-57=105$. Khớp.

Vậy $x_0+2y_0=3$.
}

\end{bt}
